\documentclass{article}
\usepackage{amsmath}
\title{Sampled Iron Ground Contact}
\author{UpstreamDrift}
\date{}
\begin{document}
\maketitle
\section{Scope and Coordinates}
Supplementary editable calculation note for issue 15523, separate from the
engineering design manual. Inputs are paired times $t_i$ in seconds and sole
heights $z_i$ in metres relative to local turf, positive upward. Ball time $t_b$
uses the same clock. Geometry, terrain sampling and ball detection are supplied
by the caller. This diagnostic applies no force, energy or spin correction.
\section{Method}
For $t\in[t_i,t_{i+1}]$, height is linearly interpolated. A crossing occurs at
$t_c=(1-f)t_i+ft_{i+1}$, where
$f=|z_i|/(|z_i|+|z_{i+1}|)$, evaluated with rescaled heights.
Let $[a_j,b_j]$ be the penetrating intervals. Then
\[
 D_{\rm pre}=\sum_j\max(0,\min(b_j,t_b)-a_j),\qquad
 d_{\max}=\max(0,-\min_i z_i).
\]
The first contact is $a_1$, or absent if no interval penetrates. The strike is
fat iff $D_{\rm pre}>0$. Tangency alone is excluded. Maximum depth includes the
whole observed record, including after the ball.
\section{Validity and Evidence}
Require finite paired samples, increasing times with finite span, at least two
points, initial height nonnegative and ball time inside the record. An ending
below turf reports only observed duration. Unobserved earlier contact and
excursions between samples cannot be inferred. No calibration coefficients or
measured physical acceptance are claimed. Analytical regression cases and exact
reproduction commands are in \texttt{IRON\_GROUND\_CONTACT.md}.
\end{document}
